\chapter{Hair Loss}

\section*{Definition}
Excessive shedding thinning hair scalp body beyond normal turnover telogen effluvium alopecia areata androgenetic alopecia distinguishing diffuse patchy pattern gradual rapid onset location extent extent severity

\section*{What Happens in Your Body}
Diverse etiations disrupting hair growth cycle anagen (growth) phase prematurely terminated pushing hairs into telogen (resting) phase causing shedding later mechanical traction damage inflammation follicle destruction hormonal influences nutritional deficiencies medications stress genetics autoimmune phenomena disrupt follicular integrity regeneration mechanisms

\section*{Causes}
\begin{itemize}
  \item Stress anemia nutritional deficiencies postpartum illness surgery weight crash dieting
  \item Medication side effects chemotherapy anticoagulants retinoids beta-blockers antidepressants
  \item Autoimmune disease alopecia areata vitiligo lupus thyroiditis
  \item Hormonal imbalance thyroid dysfunction polycystic ovary syndrome androgens male/female pattern baldness
  \item Genetic predisposition androgenetic alopecia familial hereditary pattern receding crown thinning
\end{itemize}

\section*{When to See a Doctor}
See your doctor if:
\begin{itemize}
  \item Symptoms are severe or getting worse over time
  \item Symptoms persist beyond the expected duration for this condition
  \item You have other concerning symptoms such as fever, weight loss, or bleeding
\end{itemize}

\section*{Related Symptoms}
\begin{itemize}
  \item Scalp tenderness nail changes brittle splitting clubbing hair shaft thinning inflammation redness swelling pus discharge facial hair excess facial scalp oily skin menstrual irregularities mood changes sleep disturbance appetite changes
\end{itemize}
\textbf{Important:} If this symptom persists or concerns you, your doctor will check for serious hair loss

\section*{Self-Care / Home Management}
\begin{itemize}
  \item Keep the affected area clean and dry; wash gently with mild, fragrance-free soap and lukewarm water.
  \item Avoid scratching — use cool compresses, calamine lotion, or OTC hydrocortisone cream to relieve itching.
  \item Apply moisturizer regularly to maintain skin barrier integrity, especially after bathing.
  \item Identify and avoid potential triggers (irritants, allergens, excessive heat, tight clothing).
  \item Seek medical evaluation if the rash spreads rapidly, develops blisters, shows signs of infection (pus, increasing redness, warmth), or is accompanied by fever.
\end{itemize}

\section*{How Your Doctor Will Evaluate You}
\begin{itemize}
  \item History includes onset, location, progression, associated symptoms (itching, pain, fever), exposure to triggers, and prior skin conditions.
  \item Physical examination assesses morphology (macules, papules, vesicles, etc.), distribution, and extent of skin involvement.
  \item Diagnostic tests may include skin scraping for fungal examination, patch testing for allergic contact dermatitis, or biopsy for suspicious lesions.
  \item Laboratory studies (CBC, ESR, autoimmune serologies) are obtained when systemic disease is suspected.
\end{itemize}

\section*{Treatment Options}
\begin{itemize}
  \item Topical therapies (corticosteroids, antifungals, antibiotics, emollients) are first-line for most localized skin conditions.
  \item Systemic therapy (oral antihistamines, antibiotics, antifungals, or immunosuppressants) is reserved for extensive or refractory cases.
  \item Phototherapy or biologic agents may be indicated for chronic inflammatory skin diseases (psoriasis, eczema).
  \item Referral to dermatology is appropriate for uncertain diagnoses, treatment failure, or skin lesions suspicious for malignancy.
\end{itemize}

\section*{References}

\begin{thebibliography}{99}
\bibitem{harrison-hair} Longo DL, et al. \emph{Harrison's Principles of Internal Medicine}. 21st ed. McGraw-Hill; 2022. Chapter 64: Alopecia.
\bibitem{uptodate-hair} Price VH, et al. Evaluation and diagnosis of hair loss. In: UpToDate, Post TW (Ed), UpToDate, Waltham, MA; 2025.
\bibitem{trueb} Trueb RM. Systematic approach to hair loss in women. \emph{J Dtsch Dermatol Ges}. 2023;21(4):368-382.
\bibitem{kang} Kang JS, et al. Androgenetic alopecia: an update. \emph{J Am Acad Dermatol}. 2023;89(1):S1-S8.
\bibitem{olsen} Olsen EA, et al. \emph{Disorders of Hair Growth: Diagnosis and Treatment}. 4th ed. McGraw-Hill; 2023.
\bibitem{messenger} Messenger AG, et al. Alopecia areata: a clinical review. \emph{Br J Dermatol}. 2024;190(2):175-187.
\end{thebibliography}
